% ====chapter=====
\myChapter{Mise en oeuvre et évaluation de ProOWeb}{Mise en oeuvre et évaluation}
\label{chap:mise-oeuvre-evaluation-prooweb}
\myMiniToc{mysection}{Sommaire}

\mySection{Introduction}{}

Le chapitre précédent a présenté ProOWeb comme une approche d'ingénierie permettant de passer d'un modèle BPMN enrichi à une application web métier complète et fonctionnelle. Il a décrit les modèles d'entrée, le modèle pivot, les règles de transformation et les mécanismes de préservation du code personnalisé. Le présent chapitre montre comment cette approche a été concrétisée dans un prototype et comment elle peut être évaluée à partir d'un cas organisationnel.

L'objectif n'est pas de présenter ProOWeb comme une simple démonstration technique. Dans une organisation réelle, une application web métier ne sert généralement pas un seul processus isolé. Elle regroupe plusieurs procédures, plusieurs catégories d'acteurs, des données partagées, des règles d'accès et des besoins de suivi. Pour cette raison, l'expérimentation retenue dans ce chapitre ne se limite pas à un workflow unique. Elle s'appuie sur plusieurs processus représentatifs d'une même organisation afin d'observer la capacité de ProOWeb à produire un socle applicatif cohérent.

Le chapitre est organisé en quatre parties. La première présente l'environnement de mise en oeuvre et le périmètre du prototype. La deuxième décrit l'implémentation de la chaîne de génération. La troisième présente le cas d'étude organisationnel retenu, les processus considérés et les emplacements prévus pour les modèles et captures d'écran. La dernière partie propose une évaluation de ProOWeb selon des critères liés à la complétude fonctionnelle, à la cohérence entre modèles et application, à la prise en charge de plusieurs processus et à la capacité d'évolution.

\mySection{Environnement de mise en oeuvre}{}
\label{sec:chap4-environnement}

\subsection{Objectif du prototype}

Le prototype de ProOWeb a été développé pour vérifier que l'approche proposée peut être rendue opérationnelle. Il sert à construire un projet applicatif à partir d'informations métier structurées, puis à produire les éléments nécessaires à l'exécution d'une application web. Dans cette perspective, le prototype joue deux rôles. D'une part, il constitue un éditeur et un orchestrateur de génération. D'autre part, il produit un projet cible comprenant un backend, un frontend, des fichiers de configuration et des mécanismes de suivi des fichiers générés.

Le prototype ne vise donc pas seulement la production de fragments de code. Il cherche à générer une base applicative exploitable par une équipe de développement. Cette base doit pouvoir être exécutée, testée, versionnée et complétée. Ce choix est important pour ProOWeb : l'application générée n'est pas considérée comme un résultat jetable, mais comme un point de départ structuré pour une application métier.

\subsection{Choix technologiques}

La mise en oeuvre de ProOWeb repose sur un ensemble de technologies choisies pour séparer clairement le générateur de l'application produite. Le noyau de ProOWeb est réalisé en JavaScript côté serveur. Il assure la lecture des paramètres du projet, la gestion de l'espace de travail, la préparation des modèles, la construction des contrats intermédiaires et l'appel des gabarits de génération.

L'application web générée est organisée autour d'un backend de type Spring Boot et d'un frontend de type React/Vite. Le backend porte les modules métier, les règles applicatives, les points d'accès, les mécanismes de sécurité, la persistance et les services de runtime. Le frontend expose les vues destinées aux utilisateurs : console d'accueil, tableaux de bord, listes de tâches, formulaires de traitement, vues de supervision et écrans d'administration. Des fichiers de déploiement et des scripts de vérification complètent le projet généré afin de faciliter son exécution dans plusieurs profils.

Ce choix technique intervient au chapitre 4 parce qu'il relève de la réalisation. Dans le chapitre 3, l'approche a été décrite indépendamment d'une technologie particulière. Ici, il devient nécessaire de préciser comment les concepts sont matérialisés dans le prototype. La séparation reste toutefois maintenue : ProOWeb n'est pas seulement un assemblage de frameworks ; il est d'abord une chaîne de transformation dont les technologies servent à produire une application web métier fonctionnelle.

\begin{table}[H]
\centering
\small
\caption{Principaux choix de mise en oeuvre du prototype ProOWeb}
\label{tab:chap4-choix-techniques}
\begin{tabular}{p{3.6cm}p{4.5cm}p{5.5cm}}
\hline
\textbf{Élément} & \textbf{Choix retenu} & \textbf{Rôle dans ProOWeb} \\
\hline
Noyau du générateur & JavaScript côté serveur & Lire les paramètres, piloter la génération, produire les artefacts et maintenir le manifeste. \\
Backend généré & Spring Boot & Porter la logique applicative, les modules métier, les API, la sécurité, la persistance et le runtime des processus. \\
Frontend généré & React/Vite & Fournir les écrans d'utilisation, les formulaires, la liste des tâches, les vues de suivi et les interfaces d'administration. \\
Déploiement & Fichiers Docker et profils d'exécution & Préparer l'exécution de l'application dans des environnements distincts. \\
Gouvernance & Manifeste \texttt{.prooweb-managed.json} & Identifier les fichiers générés, détecter les conflits et accompagner les régénérations. \\
\hline
\end{tabular}
\end{table}

\subsection{Organisation générale du prototype}

Le prototype est organisé en deux grands ensembles. Le premier correspond au générateur ProOWeb. Il contient les services de gestion du workspace, les services de modélisation de processus, les contrats de génération, les gabarits et les routes permettant de piloter l'outil. Le second correspond au projet applicatif généré. Il contient le backend, le frontend, les fichiers de déploiement et les scripts de construction ou de vérification.

Cette organisation traduit une séparation volontaire. Le générateur conserve la responsabilité de produire et de suivre les fichiers. L'application générée conserve la responsabilité d'exécuter les fonctions métier. Une telle séparation permet d'éviter que le prototype devienne un simple environnement fermé. L'organisation visée donne au développeur la possibilité d'ouvrir le code généré, de le comprendre et de le faire évoluer.

\begin{figure}[H]
\centering
\fbox{\parbox[c][5cm][c]{0.9\textwidth}{\centering
Emplacement réservé à une figure présentant l'organisation générale du prototype ProOWeb : générateur, modèles, contrats, templates et application produite.
}}
\caption{Organisation générale du prototype ProOWeb}
\label{fig:chap4-organisation-prototype}
\end{figure}

\mySection{Implémentation de la chaîne de génération}{}
\label{sec:chap4-implementation-generation}

\subsection{Initialisation du projet cible}

La génération commence par l'initialisation du projet cible. Cette étape permet de définir les informations de base de l'application à produire : nom du projet, espace de travail, package de base, options d'authentification, activation de certains modules, profils de déploiement et paramètres de configuration. Ces informations ne sont pas liées à un seul processus ; elles décrivent le cadre général de l'application métier que l'organisation veut obtenir.

Dans ProOWeb, cette initialisation est importante parce qu'elle fixe l'identité de l'application générée. Elle permet d'éviter une génération anonyme ou difficile à maintenir. Le projet produit possède une structure, une configuration et des conventions qui pourront être reprises par les développeurs. L'application générée peut ainsi être intégrée dans un cycle de développement classique.

\begin{figure}[H]
\centering
\fbox{\parbox[c][5cm][c]{0.9\textwidth}{\centering
Emplacement réservé à une capture du wizard ou de l'écran d'initialisation du projet ProOWeb.
}}
\caption{Écran d'initialisation du projet dans ProOWeb}
\label{fig:chap4-initialisation-projet}
\end{figure}

\subsection{Gestion des modèles de processus}

Après l'initialisation, ProOWeb prend en charge les modèles de processus. Chaque processus est décrit par un modèle BPMN et complété par une spécification \textit{spec-v1}. Le BPMN garde le rôle de support de workflow. La spécification enrichie précise les éléments nécessaires à l'application : rôles candidats, règles d'affectation, entrées, sorties, données partagées, règles de visibilité et stratégies de stockage.

Le prototype gère également le cycle de vie des versions de processus. Un modèle peut être en brouillon, validé, déployé ou retiré. Cette logique correspond à une réalité organisationnelle : une procédure peut être modifiée avant son entrée en production ; une ancienne version peut rester consultable ; une nouvelle version peut être déployée lorsque les règles métier changent. La gestion des versions permet donc de ne pas confondre la modélisation, la validation et l'exploitation.

\begin{figure}[H]
\centering
\fbox{\parbox[c][5cm][c]{0.9\textwidth}{\centering
Emplacement réservé à une capture de l'espace de gestion des modèles de processus et de leurs versions.
}}
\caption{Gestion des modèles et des versions de processus}
\label{fig:chap4-gestion-modeles-processus}
\end{figure}

\subsection{Validation et construction des contrats}

La validation constitue une étape centrale de l'implémentation. Le prototype vérifie la cohérence entre le modèle BPMN et la spécification \textit{spec-v1}. Une activité enrichie doit exister dans le modèle de processus. Une tâche manuelle doit être associée à des rôles capables de la traiter. Un champ utilisé dans un formulaire doit provenir d'une source connue. Une donnée partagée doit posséder une définition suffisante pour être projetée dans l'application.

Lorsque les informations sont cohérentes, ProOWeb construit des contrats intermédiaires. Ces contrats ne sont pas encore du code source ; ils représentent une forme stabilisée du futur comportement applicatif. Le \textit{runtime contract} décrit les processus, les activités, les transitions et les règles d'exécution. Le \textit{data contract} décrit les données manipulées et leur stockage. Les catalogues de formulaires et d'interfaces préparent les vues qui seront exposées aux utilisateurs.

\begin{table}[H]
\centering
\small
\caption{Contrats produits avant la génération des fichiers}
\label{tab:chap4-contrats-generation}
\begin{tabular}{p{3.8cm}p{5.0cm}p{4.8cm}}
\hline
\textbf{Contrat} & \textbf{Contenu principal} & \textbf{Utilisation dans l'application générée} \\
\hline
\textit{Runtime contract} & Processus, activités, transitions, statuts et règles d'exécution & Construire le moteur d'exécution, les API de tâches et les vues de suivi. \\
\textit{Data contract} & Données partagées, champs, sources et mappings & Générer la persistance, les projections de données et les champs de formulaires. \\
Catalogue des formulaires & Activités manuelles, champs affichés, champs produits et contraintes & Générer les formulaires de traitement associés aux tâches. \\
Catalogue d'interface & Vues, listes de tâches, actions et règles d'affichage & Générer les composants frontend liés au processus. \\
Registre des processus & Processus déployés, versions, métadonnées et visibilité & Initialiser le runtime et organiser l'accès aux processus disponibles. \\
\hline
\end{tabular}
\end{table}

\subsection{Génération du backend}

La génération du backend produit la partie applicative de l'application web métier. Elle comprend les modules de domaine, les cas d'utilisation, les adaptateurs techniques, les points d'accès et les éléments de configuration. La logique générée vise notamment à démarrer une instance de processus, créer ou afficher les tâches, compléter une activité, appliquer les transitions, filtrer les données visibles et enregistrer l'historique des actions.

Le backend généré suit une organisation modulaire. Cette organisation aide à séparer les responsabilités : identité, organisation, notifications, sécurité, runtime de processus et données métier. Elle facilite aussi la lecture du code généré. Dans un mémoire orienté recherche appliquée, ce point est essentiel : le résultat de ProOWeb doit être compréhensible par un développeur, pas seulement exécutable par une machine.

\subsection{Génération du frontend}

La génération du frontend produit les interfaces nécessaires aux utilisateurs de l'organisation. Les vues générées ne se limitent pas à un formulaire isolé. Elles couvrent la consultation des processus disponibles, la liste des tâches à traiter, les tâches à affecter, les tâches visibles, les instances suivies, la complétion des activités et la supervision.

Cette partie rend visible l'objectif de ProOWeb. Un acteur métier ne manipule pas directement le modèle BPMN ni le fichier \textit{spec-v1}. Il utilise une application. Le frontend généré est donc le point de rencontre entre l'approche de modélisation et l'usage réel. Si cette interface ne présente pas correctement les tâches, les rôles et les données, l'application générée reste incomplète.

\begin{figure}[H]
\centering
\fbox{\parbox[c][5cm][c]{0.9\textwidth}{\centering
Emplacement réservé à une capture de la liste des tâches ou du tableau de bord runtime généré.
}}
\caption{Vue générée de traitement des tâches}
\label{fig:chap4-vue-taches-generee}
\end{figure}

\subsection{Génération des éléments de déploiement et de vérification}

Le prototype génère également des éléments qui facilitent l'exécution et la vérification de l'application. Les profils d'exécution permettent de distinguer les contextes de développement, de démonstration, de test ou de production. Les scripts de construction et de vérification donnent un point d'entrée simple pour compiler, tester ou démarrer l'application produite.

Cette dimension peut paraître secondaire, mais elle participe à la notion d'application fonctionnelle. Une application métier générée ne doit pas seulement contenir des classes et des composants. Elle doit pouvoir être lancée, vérifiée et livrée. Les éléments de déploiement contribuent donc à rendre le résultat utilisable par une équipe technique.

\subsection{Préservation et régénération contrôlée}

Le prototype met en place un manifeste des fichiers générés. Ce manifeste garde la trace des fichiers produits, de leurs catégories et de leurs empreintes. Lorsqu'une nouvelle génération est lancée, ProOWeb peut comparer l'état attendu avec l'état présent dans le projet cible. Si un fichier généré a été modifié manuellement, l'outil peut signaler le conflit, créer une sauvegarde ou appliquer une stratégie de migration.

Ce mécanisme est indispensable dans le cadre d'une application destinée à une organisation. Une fois l'application générée, des adaptations locales peuvent apparaître : ajustement d'un libellé, ajout d'une règle particulière, intégration à un service externe ou modification d'une vue. La régénération ne doit pas effacer ces adaptations sans contrôle. ProOWeb cherche donc à préserver le lien entre modèle et code, tout en respectant l'intervention humaine.

\mySection{Cas d'étude organisationnel}{}
\label{sec:chap4-cas-etude}

\subsection{Présentation de l'organisation retenue}

Le cas d'étude retenu représente une organisation administrative qui souhaite numériser plusieurs procédures internes. Cette organisation reçoit des demandes, les affecte à des agents, les fait valider par des responsables, conserve les données produites et doit permettre à certains acteurs de suivre l'évolution des dossiers. Elle correspond volontairement à un contexte courant : plusieurs services travaillent sur des procédures distinctes, mais avec des acteurs, des informations et des besoins de supervision partagés.

Le choix d'un cas organisationnel répond à une limite fréquente des démonstrations de génération d'applications : elles se concentrent sur un seul processus et montrent que la chaîne de transformation fonctionne dans un cas réduit. Or, une application web métier utilisée par une organisation doit regrouper plusieurs processus. Elle doit offrir une expérience cohérente, partager certaines données, appliquer des règles de sécurité homogènes et permettre aux superviseurs d'avoir une vue globale.

\begin{figure}[H]
\centering
\fbox{\parbox[c][5cm][c]{0.9\textwidth}{\centering
Emplacement réservé à un schéma présentant l'organisation étudiée, ses services et les acteurs impliqués.
}}
\caption{Vue générale de l'organisation retenue pour le cas d'étude}
\label{fig:chap4-organisation-cas-etude}
\end{figure}

\subsection{Acteurs et responsabilités}

Les acteurs retenus dans le cas d'étude sont volontairement simples, afin de garder l'expérimentation lisible tout en couvrant les responsabilités essentielles d'une application métier. Le demandeur initie une procédure et fournit les informations nécessaires. L'agent traitant vérifie les données, complète le dossier et prépare une décision. Le validateur approuve ou rejette la demande selon les règles définies. Le superviseur suit l'évolution des instances et intervient lorsque des tâches restent bloquées. L'administrateur configure les paramètres de base et gère les comptes ou les rôles.

\begin{table}[H]
\centering
\small
\caption{Acteurs du cas d'étude organisationnel}
\label{tab:chap4-acteurs-cas-etude}
\begin{tabular}{p{3.2cm}p{5.3cm}p{5.0cm}}
\hline
\textbf{Acteur} & \textbf{Responsabilité principale} & \textbf{Interaction avec l'application générée} \\
\hline
Demandeur & Soumettre une demande et consulter son avancement & Remplit un formulaire initial et suit ses dossiers. \\
Agent traitant & Vérifier les informations et compléter le traitement & Reçoit les tâches affectées et renseigne les décisions intermédiaires. \\
Validateur & Approuver, rejeter ou retourner un dossier & Consulte les données utiles et exécute les actions de validation. \\
Superviseur & Suivre les instances et résoudre les blocages & Observe les tâches, les retards, les affectations et les événements. \\
Administrateur & Configurer les rôles et les paramètres organisationnels & Gère les utilisateurs, les unités, les notifications et les options. \\
\hline
\end{tabular}
\end{table}

\subsection{Cartographie des processus retenus}

Trois processus ont été retenus pour représenter le fonctionnement de l'organisation. Le premier concerne la demande administrative simple. Il couvre la soumission d'une demande, sa vérification, sa validation et sa clôture. Le deuxième concerne le traitement d'une réclamation. Il met en évidence la réception d'un signalement, son analyse, son affectation et la réponse au demandeur. Le troisième concerne une demande d'approvisionnement interne. Il introduit une logique de validation hiérarchique et de suivi plus marquée.

Ces processus ne sont pas identiques, mais ils partagent plusieurs éléments : des acteurs communs, des données de demande, des règles d'affectation, des statuts d'instance, des notifications et des besoins de supervision. Cette situation permet d'évaluer la capacité de ProOWeb à produire une application cohérente pour plusieurs workflows appartenant au même système métier.

\begin{table}[H]
\centering
\small
\caption{Processus métier retenus dans le cas d'étude}
\label{tab:chap4-processus-retenus}
\begin{tabular}{p{3.2cm}p{4.4cm}p{3.0cm}p{3.4cm}}
\hline
\textbf{Processus} & \textbf{Objectif métier} & \textbf{Acteurs clés} & \textbf{Résultat attendu} \\
\hline
P1 : demande administrative & Recevoir, vérifier et valider une demande formulée par un usager interne & Demandeur, agent, validateur & Demande acceptée, rejetée ou retournée. \\
P2 : réclamation & Traiter une réclamation et produire une réponse suivie & Demandeur, agent, superviseur & Réclamation clôturée avec réponse. \\
P3 : approvisionnement interne & Demander et valider un besoin de matériel ou de service & Demandeur, validateur, superviseur & Demande approuvée ou refusée. \\
\hline
\end{tabular}
\end{table}

\begin{figure}[H]
\centering
\fbox{\parbox[c][5cm][c]{0.9\textwidth}{\centering
Emplacement réservé à une cartographie montrant les trois processus retenus et leurs liens avec les acteurs de l'organisation.
}}
\caption{Cartographie des processus métier retenus}
\label{fig:chap4-cartographie-processus}
\end{figure}

\subsection{Description du processus P1 : demande administrative}

Le processus de demande administrative commence lorsqu'un demandeur soumet un formulaire contenant l'objet de la demande, une description et les pièces nécessaires. L'application crée une instance de processus et rend la tâche visible à un agent traitant. L'agent vérifie la complétude des informations. Si le dossier est incomplet, il peut être retourné au demandeur. Si le dossier est complet, il est transmis à un validateur. Le validateur prend une décision finale : accepter, rejeter ou demander une correction.

Ce processus illustre un cas classique de workflow à validation. Il permet de vérifier que ProOWeb génère les formulaires, les transitions, les règles de visibilité et les statuts nécessaires. Il permet également d'observer le traitement des retours, car une demande ne suit pas toujours un chemin linéaire.

\begin{figure}[H]
\centering
\fbox{\parbox[c][5cm][c]{0.9\textwidth}{\centering
Emplacement réservé au modèle BPMN du processus P1 : demande administrative.
}}
\caption{Modèle BPMN du processus de demande administrative}
\label{fig:chap4-bpmn-demande-administrative}
\end{figure}

\subsection{Description du processus P2 : traitement d'une réclamation}

Le processus de réclamation commence par la saisie d'un problème rencontré par un usager interne. La réclamation est ensuite qualifiée par un agent. Selon sa nature, elle peut être traitée directement ou transmise à un responsable. Lorsque la réponse est rédigée, elle est renvoyée au demandeur et l'instance est clôturée. Si la réclamation révèle un problème plus large, le superviseur peut consulter l'historique et suivre les actions entreprises.

Ce processus permet d'évaluer des besoins légèrement différents de ceux du premier. Il implique plus fortement le suivi, la justification des décisions et la consultation de l'historique. Il montre aussi que l'application générée doit pouvoir gérer des processus dont la valeur métier se trouve autant dans la trace que dans la décision finale.

\begin{figure}[H]
\centering
\fbox{\parbox[c][5cm][c]{0.9\textwidth}{\centering
Emplacement réservé au modèle BPMN du processus P2 : traitement d'une réclamation.
}}
\caption{Modèle BPMN du processus de traitement d'une réclamation}
\label{fig:chap4-bpmn-reclamation}
\end{figure}

\subsection{Description du processus P3 : approvisionnement interne}

Le processus d'approvisionnement interne permet à un service de demander un bien ou un service nécessaire à son fonctionnement. Le demandeur précise le besoin, la justification, la quantité et l'urgence. La demande est ensuite examinée par un validateur. Selon le montant ou la nature du besoin, elle peut être transmise à un superviseur. L'instance se termine par une approbation, un rejet ou une demande de complément.

Ce troisième processus introduit une logique de décision plus hiérarchique. Il permet de tester les stratégies d'affectation, la visibilité par rôle et le suivi des tâches non traitées. Dans une application organisationnelle, ce type de processus est important parce qu'il combine données, responsabilité, contrôle et traçabilité.

\begin{figure}[H]
\centering
\fbox{\parbox[c][5cm][c]{0.9\textwidth}{\centering
Emplacement réservé au modèle BPMN du processus P3 : approvisionnement interne.
}}
\caption{Modèle BPMN du processus d'approvisionnement interne}
\label{fig:chap4-bpmn-approvisionnement}
\end{figure}

\subsection{Spécifications enrichies associées}

Pour chaque processus, le modèle BPMN est complété par une spécification \textit{spec-v1}. Cette spécification précise les informations qui ne sont pas naturellement portées par le diagramme : rôles autorisés, données d'entrée, sorties produites, règles de stockage et règles de visibilité. Elle permet à ProOWeb de passer d'un workflow à une application métier utilisable.

L'extrait suivant illustre la forme attendue d'une activité enrichie. Il ne constitue pas la totalité de la spécification ; il montre seulement les informations nécessaires à la génération d'une tâche manuelle.

\begin{lstlisting}[caption={Exemple simplifié d'activité enrichie pour le cas d'étude},label={lst:chap4-exemple-spec}]
{
  "activityId": "verifier_dossier",
  "type": "manual",
  "candidateRoles": ["AGENT_TRAITANT"],
  "inputs": ["demande.objet", "demande.description", "demande.pieces"],
  "outputs": ["decisionVerification", "commentaireAgent"],
  "visibility": {
    "activityViewerRoles": ["AGENT_TRAITANT", "SUPERVISEUR"],
    "dataViewerRoles": ["VALIDATEUR"]
  },
  "storage": "BOTH"
}
\end{lstlisting}

\begin{table}[H]
\centering
\small
\caption{Synthèse des spécifications enrichies par processus}
\label{tab:chap4-synthese-specs}
\begin{tabular}{p{2.5cm}p{3.5cm}p{3.0cm}p{4.6cm}}
\hline
\textbf{Processus} & \textbf{Rôles principaux} & \textbf{Données partagées} & \textbf{Écrans attendus} \\
\hline
P1 & Demandeur, agent, validateur & Demande, décision, pièces & Formulaire de demande, vérification, validation, suivi. \\
P2 & Demandeur, agent, superviseur & Réclamation, réponse, historique & Dépôt de réclamation, qualification, réponse, supervision. \\
P3 & Demandeur, validateur, superviseur & Besoin, justification, avis & Demande d'achat, validation, arbitrage, clôture. \\
\hline
\end{tabular}
\end{table}

\subsection{Application web métier générée}

À partir des trois processus et de leurs spécifications, ProOWeb doit produire une application web métier unique. Cette application regroupe les processus disponibles, affiche les tâches selon les rôles, présente les formulaires adaptés à chaque activité et conserve les données utiles à la progression des instances. L'utilisateur ne voit pas trois applications séparées ; il utilise un même environnement dans lequel plusieurs procédures sont accessibles.

Cette unité applicative constitue un résultat important. Elle montre que ProOWeb peut soutenir une logique organisationnelle. Les processus restent distincts du point de vue métier, mais ils partagent une infrastructure d'exécution, une gestion des rôles, des mécanismes de suivi et une expérience utilisateur commune.

\begin{figure}[H]
\centering
\fbox{\parbox[c][5cm][c]{0.9\textwidth}{\centering
Emplacement réservé à une capture de la page d'accueil ou du tableau de bord de l'application générée.
}}
\caption{Tableau de bord de l'application web métier générée}
\label{fig:chap4-dashboard-genere}
\end{figure}

\begin{figure}[H]
\centering
\fbox{\parbox[c][5cm][c]{0.9\textwidth}{\centering
Emplacement réservé à une capture d'un formulaire généré pour l'une des activités manuelles.
}}
\caption{Exemple de formulaire généré pour une activité métier}
\label{fig:chap4-formulaire-genere}
\end{figure}

\begin{figure}[H]
\centering
\fbox{\parbox[c][5cm][c]{0.9\textwidth}{\centering
Emplacement réservé à une capture d'une vue de supervision ou d'historique des instances.
}}
\caption{Vue de suivi et de supervision des instances}
\label{fig:chap4-supervision-generee}
\end{figure}

\mySection{Évaluation de ProOWeb}{}
\label{sec:chap4-evaluation}

\subsection{Méthode d'évaluation}

L'évaluation menée dans ce mémoire est une évaluation fonctionnelle et qualitative du prototype. Elle cherche à vérifier si ProOWeb permet effectivement de générer une application web métier exploitable à partir de plusieurs processus d'une même organisation. Elle ne prétend pas mesurer définitivement la productivité dans tous les contextes, ni comparer ProOWeb à l'ensemble des plateformes existantes. Elle se concentre sur les résultats obtenus dans le cadre du prototype.

La méthode retenue consiste à observer les éléments produits par la chaîne de génération et à les comparer aux besoins exprimés dans les modèles d'entrée. Les critères portent sur la couverture fonctionnelle, la cohérence entre BPMN et application, la prise en compte des rôles, la génération des formulaires, la gestion des données, la supervision, la capacité multi-processus et la régénération contrôlée.

\begin{table}[H]
\centering
\small
\caption{Critères retenus pour l'évaluation}
\label{tab:chap4-criteres-evaluation}
\begin{tabular}{p{3.8cm}p{6.0cm}p{3.8cm}}
\hline
\textbf{Critère} & \textbf{Question posée} & \textbf{Mode d'observation} \\
\hline
Complétude fonctionnelle & Les éléments nécessaires à l'exécution de l'application sont-ils générés ? & Analyse des modules produits et exécution du parcours. \\
Cohérence BPMN-application & Les activités et transitions du modèle sont-elles représentées dans l'application ? & Comparaison entre modèle, contrats et vues générées. \\
Gestion des rôles & Les acteurs voient-ils uniquement les tâches et données prévues ? & Observation des listes de tâches et règles d'accès. \\
Gestion des données & Les entrées, sorties et données partagées sont-elles prises en compte ? & Vérification des formulaires, mappings et projections. \\
Capacité multi-processus & L'application regroupe-t-elle plusieurs processus dans un même environnement ? & Observation du registre, des tableaux de bord et des instances. \\
Évolution contrôlée & Une régénération peut-elle être suivie et sécurisée ? & Analyse du manifeste, des sauvegardes et rapports de migration. \\
\hline
\end{tabular}
\end{table}

\subsection{Résultats observés}

Les résultats attendus se situent à deux niveaux. Au niveau du générateur, ProOWeb doit lire les modèles, valider les informations, produire les contrats et écrire les artefacts du projet cible. Au niveau de l'application générée, les utilisateurs doivent pouvoir accéder aux processus, démarrer une instance, traiter les tâches, consulter les données autorisées et suivre l'état des dossiers.

Dans le cas d'étude, les trois processus retenus permettent d'observer ces comportements dans des situations différentes. Le processus de demande administrative met l'accent sur la validation. Le processus de réclamation met l'accent sur le suivi et la réponse. Le processus d'approvisionnement met l'accent sur la décision hiérarchique. Leur génération dans une même application permet de vérifier que ProOWeb ne dépend pas d'un seul workflow figé.

\begin{table}[H]
\centering
\small
\caption{Synthèse des résultats attendus sur le cas d'étude}
\label{tab:chap4-resultats-attendus}
\begin{tabular}{p{4.0cm}p{4.8cm}p{4.8cm}}
\hline
\textbf{Élément évalué} & \textbf{Résultat observé ou attendu} & \textbf{Interprétation} \\
\hline
Registre des processus & Les trois processus sont disponibles dans la même application & ProOWeb soutient une logique organisationnelle et non un processus isolé. \\
Listes de tâches & Les tâches sont filtrées selon les rôles & Les règles déclarées dans \textit{spec-v1} sont projetées vers l'application. \\
Formulaires & Les champs affichés correspondent aux activités & Les données d'entrée et de sortie sont exploitées par la génération. \\
Runtime & Les instances progressent selon les transitions prévues & Le modèle BPMN reste relié au comportement applicatif. \\
Supervision & Les instances et actions peuvent être suivies & L'application générée répond aux besoins de contrôle organisationnel. \\
Manifeste & Les fichiers générés sont suivis & La régénération peut être analysée et maîtrisée. \\
\hline
\end{tabular}
\end{table}

\subsection{Analyse de la capacité multi-processus}

La capacité multi-processus constitue l'apport principal de l'expérimentation. Elle montre que ProOWeb ne doit pas être réduit à un générateur de formulaires pour une procédure isolée. Une organisation a besoin d'une application où plusieurs processus coexistent. Cette coexistence exige une gestion commune des utilisateurs, des rôles, des données, des permissions et des vues de suivi.

Dans le cas d'étude, les processus P1, P2 et P3 partagent certains acteurs tout en conservant des objectifs différents. Cette configuration permet de vérifier que le modèle de génération ne confond pas les workflows. Chaque processus possède ses activités, ses transitions et ses formulaires. En même temps, l'application générée conserve un socle commun : authentification, rôles, tableau de bord, listes de tâches, historique et supervision.

Cette séparation entre indépendance des processus et unité de l'application est essentielle. Elle correspond à l'usage réel d'un système d'information métier. Les services d'une organisation ne veulent pas nécessairement trois applications différentes pour trois procédures. Ils veulent un système cohérent dans lequel les processus sont identifiables, mais intégrés.

\subsection{Apports de l'approche}

L'expérimentation met en évidence plusieurs apports. Le premier est la continuité entre la modélisation et l'application produite. Les activités BPMN ne restent pas seulement des éléments graphiques ; elles deviennent des tâches, des vues, des actions et des transitions. Le deuxième apport est la prise en compte des informations applicatives par \textit{spec-v1}. Les rôles, les données et les règles de visibilité ne sont pas ajoutés après coup dans le code ; ils font partie de la source de génération.

Le troisième apport concerne la structuration du projet généré. ProOWeb produit une application organisée, avec un backend, un frontend, des fichiers de configuration, des profils d'exécution et un mécanisme de suivi des fichiers. Cette organisation rend le résultat plus proche d'un projet logiciel exploitable. Le quatrième apport concerne la régénération contrôlée. En suivant les fichiers générés, ProOWeb prépare l'évolution de l'application au lieu de supposer que le processus ne changera plus.

\subsection{Limites observées}

Le prototype présente également des limites. La première concerne la qualité des modèles d'entrée. Si le BPMN est ambigu ou si la spécification \textit{spec-v1} est incomplète, la génération ne peut pas produire une application correcte. ProOWeb peut signaler des incohérences, mais il ne peut pas remplacer l'analyse métier. La deuxième limite concerne la couverture des comportements très spécifiques. Certaines règles complexes peuvent nécessiter une extension manuelle ou un point d'intégration particulier.

La troisième limite concerne l'évaluation elle-même. Le cas organisationnel retenu permet de vérifier la faisabilité et la cohérence de l'approche, mais il ne suffit pas à conclure sur tous les domaines métier possibles. Une validation plus large demanderait d'appliquer ProOWeb à d'autres organisations, avec des processus plus nombreux, des données plus complexes et des contraintes de sécurité plus fortes.

Enfin, l'application générée reste un socle. Elle peut être fonctionnelle et exploitable, mais elle n'élimine pas le travail du développeur. Dans un projet réel, une équipe devra encore adapter certains écrans, connecter des services externes, ajuster des règles spécifiques et enrichir les tests. L'intérêt de ProOWeb est de réduire la distance entre le modèle métier et cette base applicative, non de supprimer toute intervention humaine.

\subsection{Discussion}

Les résultats obtenus permettent de situer ProOWeb comme une approche d'ingénierie orientée génération d'applications web métier. Sa valeur ne réside pas uniquement dans la rapidité de production du code, mais dans la cohérence qu'il maintient entre processus, données, rôles et application. Cette cohérence est particulièrement importante lorsque plusieurs processus doivent être réunis dans une même application.

Le cas d'étude montre aussi que la spécification enrichie joue un rôle déterminant. BPMN décrit le déroulement du travail, mais l'application nécessite des informations supplémentaires. En rendant ces informations explicites, ProOWeb rend la génération plus contrôlable. Le prix à payer est un effort de spécification plus important au départ. Cet effort paraît néanmoins justifié lorsque l'objectif est de produire une application métier complète plutôt qu'un prototype visuel.

\mySection{Conclusion}{}

Ce chapitre a présenté la mise en oeuvre et l'évaluation de ProOWeb. Le prototype développé matérialise les principes définis au chapitre précédent : lecture des modèles BPMN, enrichissement par \textit{spec-v1}, validation, construction de contrats intermédiaires, génération du backend, génération du frontend, production des éléments de déploiement et suivi des fichiers générés.

L'expérimentation a été pensée à l'échelle d'une organisation et non d'un processus isolé. Trois processus métier ont été retenus afin de montrer que ProOWeb peut produire une application web métier unique, capable de gérer plusieurs workflows, plusieurs acteurs, des données partagées et des besoins de supervision. Cette orientation renforce l'intérêt de l'approche, car elle correspond mieux à la manière dont les organisations utilisent réellement les applications métier.

L'évaluation montre que ProOWeb apporte une réponse structurée au passage du modèle de processus vers l'application fonctionnelle. Elle met aussi en évidence les limites du prototype : dépendance à la qualité des modèles, nécessité d'une spécification enrichie, couverture partielle de certains cas très spécifiques et besoin d'expérimentations plus larges. Ces limites ouvrent naturellement sur la conclusion générale, où seront rappelées les contributions du mémoire et les perspectives d'amélioration de ProOWeb.
